\documentclass[11pt]{article}

\usepackage[margin=1.1in]{geometry}
\usepackage{amsmath,amssymb,amsthm,mathtools}
\usepackage{booktabs}
\usepackage{enumitem}
\usepackage[colorlinks=true,linkcolor=blue,citecolor=blue,urlcolor=blue]{hyperref}

% ---------- theorem environments ----------
\theoremstyle{plain}
\newtheorem{theorem}{Theorem}[section]
\newtheorem{proposition}[theorem]{Proposition}
\newtheorem{lemma}[theorem]{Lemma}
\newtheorem{corollary}[theorem]{Corollary}
\newtheorem{fact}[theorem]{Fact}            % imported results
\newtheorem{principle}[theorem]{Principle}
\theoremstyle{definition}
\newtheorem{definition}[theorem]{Definition}
\newtheorem{assumption}[theorem]{Assumption}
\newtheorem{prediction}[theorem]{Prediction}
\theoremstyle{remark}
\newtheorem{remark}[theorem]{Remark}

% ---------- macros ----------
\newcommand{\E}{\mathbb{E}}
\newcommand{\Prob}{\mathbb{P}}
\newcommand{\R}{\mathbb{R}}
\newcommand{\cC}{\mathcal{C}}
\newcommand{\cD}{\mathcal{D}}
\newcommand{\cF}{\mathcal{F}}
\newcommand{\cL}{\mathcal{L}}
\newcommand{\cT}{\mathcal{T}}
\newcommand{\cY}{\mathcal{Y}}
\newcommand{\cX}{\mathcal{X}}
\newcommand{\cR}{\mathcal{R}}
\newcommand{\post}{\Pi}
\DeclareMathOperator*{\argmax}{arg\,max}

\title{Prompting as Training:\\
\large Context Optimization as Learning in a Semi-Parametric Predictive Model}
\author{Working notes --- agent-families project}
\date{June 2026 --- draft 0}

\begin{document}
\maketitle

\begin{abstract}
We reframe the agent-families system's learning mechanism --- accumulating a
governed library of textual insights injected into context --- as
\emph{model training in a different coordinate system}. The foundational
frame is Bayesian (\S\ref{sec:frame}): a prompt is a conditioning operation
on a predictive distribution, a weight update is a reparameterization of the
same distribution, and in an idealized exchangeable predictive model the two
operators \emph{coincide} on predictives (the Equivalence Principle,
Proposition~\ref{prop:definetti}). This idealization is supported from below
by exact dual-form theorems for linear attention
(\S\ref{sec:dual}: a trained linear layer \emph{is} attention over its
training data, Theorem~\ref{thm:dualform}) and bounded from above by
elicitation-ceiling results (\S\ref{sec:ceiling}: context can only elicit
what the frozen model represents). Within this frame the skill library is
the trainable nonparametric component of a semi-parametric model, and every
mechanism of the system acquires a training-theoretic dual
(\S\ref{sec:dictionary}). The reframe yields falsifiable predictions
(\S\ref{sec:predictions}) and a precise answer to the question ``is this
just a poor man's SGD?'' (\S\ref{sec:regime}: no --- it is training in a
\emph{regime} of few, large, semantically generated, individually
hypothesis-tested updates, where the dual representation dominates).
Earlier results --- the measurement-theoretic outer loop of
\texttt{theory.tex} --- remain valid and are catalogued as companion modules
to revisit (\S\ref{sec:revisit}).
\end{abstract}

\tableofcontents

%======================================================================
\section{Thesis and relation to prior work in this repository}
%======================================================================

\paragraph{Thesis.} The agent-families system does not merely \emph{use} an
LLM with helpful context; it \emph{trains a model}. The trainable parameters
are not weights but the contents of a curated text library $L$, consulted by
retrieval at inference time. The claim made precise in this document:
\begin{quote}\itshape
Conditioning a predictive model on curated context and updating its
parameters on curated data are two coordinate representations of the same
learning operation. The system performs the former; consequently the
mathematics of the latter --- learning curves, capacity control,
regularization, forgetting, curriculum, distillation --- transfers, with the
translation dictionary of \S\ref{sec:dictionary} and the regime caveats of
\S\ref{sec:regime}.
\end{quote}

\paragraph{Relation to \texttt{theory.tex}.} The companion document
formalizes the \emph{outer loop}: a budgeted zeroth-order optimization with
anytime-valid gates, drift convergence, submodular governance, and an
adversarially maintained instrument. Nothing there is displaced. This
document supplies what the outer loop treated as opaque: the
\emph{semantics of the update operator} --- what kind of object a promoted
insight is, and why moves on $L$ are training steps in a rigorous sense.
Pointers to the retained modules are collected in \S\ref{sec:revisit}.

\paragraph{Status conventions.} As in the companion document:
\emph{proved here} (self-contained proofs), \emph{imported}
(\textbf{Fact} environments with citations), \emph{assumed}
(\textbf{Assumption} environments; empirical, estimable).

%======================================================================
\section{The foundational frame: one predictive distribution, two update
operators}\label{sec:frame}
%======================================================================

\subsection{Setup}

\begin{definition}[Predictive model; the two operators]\label{def:operators}
Let $p(\,y \mid x,\, C\,;\, \theta\,)$ be a conditional predictive
distribution over outputs $y \in \cY$ given an input $x \in \cX$ and a
context $C \in \cC$ (a finite token sequence), with parameters $\theta$.
Two update operators act on the induced predictive:
\[
  \underbrace{T_{\mathrm{ctx}}[D]\,:\; p(\cdot \mid x) \;\longmapsto\;
  p(\cdot \mid x,\, D)}_{\text{conditioning (prompting / retrieval)}}
  \qquad\qquad
  \underbrace{T_{\mathrm{wt}}[D]\,:\; \theta \;\longmapsto\;
  \theta'(D)}_{\text{parameter update (training)}}
\]
where $D$ is data (demonstrations, instructions, insights). The system under
study applies only $T_{\mathrm{ctx}}$, with $D$ drawn from a library $L$ by a
retrieval policy $R$.
\end{definition}

\subsection{The Equivalence Principle}

The cleanest general statement lives in the Bayesian idealization. Recall de
Finetti's representation theorem: an infinitely exchangeable sequence
$(Y_1, Y_2, \dots)$ admits a representation as a mixture of i.i.d.\ models,
\[
  p(y_1, \dots, y_n) \;=\; \int \prod_{i=1}^{n} p(y_i \mid \vartheta)\;
  d\post(\vartheta),
\]
for a unique mixing measure $\post$ over a latent parameter $\vartheta$
\cite{definetti,bernardosmith}.

\begin{proposition}[Equivalence Principle: conditioning $=$ training, on
predictives]\label{prop:definetti}
Let the data-generating process for a task family be exchangeable with de
Finetti measure $\post$ over a latent task/parameter $\vartheta$. For a
dataset $D = (d_1, \dots, d_n)$ define:
\begin{enumerate}[label=(\roman*),leftmargin=2.5em]
  \item the \emph{context-updated} predictive
        $p_{\mathrm{ctx}}(y \mid x) = p(y \mid x, D)$, obtained by
        conditioning the joint exchangeable model on $D$;
  \item the \emph{parameter-updated} predictive
        $p_{\mathrm{wt}}(y \mid x) = \int p(y \mid x, \vartheta)\,
        d\post(\vartheta \mid D)$, obtained by replacing the prior $\post$
        with the posterior $\post(\cdot \mid D)$ and integrating it out.
\end{enumerate}
Then $p_{\mathrm{ctx}} = p_{\mathrm{wt}}$ identically: \emph{learning by
conditioning and learning by parameter update induce the same predictive
distribution.}
\end{proposition}

\begin{proof}
By the representation theorem, conditioning the mixture on $D$ yields
\[
  p(y \mid x, D)
  \;=\; \frac{\int p(y \mid x, \vartheta) \prod_i p(d_i \mid \vartheta)\,
  d\post(\vartheta)}{\int \prod_i p(d_i \mid \vartheta)\, d\post(\vartheta)}
  \;=\; \int p(y \mid x, \vartheta)\; d\post(\vartheta \mid D),
\]
where the second equality is Bayes' rule applied to the mixing measure. The
right side is precisely $p_{\mathrm{wt}}(y \mid x)$.
\end{proof}

\begin{remark}[Why this is the right foundation]
Proposition~\ref{prop:definetti} is elementary, but it is the load-bearing
reframe: \emph{prompting and training are not analogous; in the exchangeable
ideal they are equal as operators on predictives}. Everything else in this
document manages the two directions of departure from the ideal: real
attention mechanisms only approximate the conditioning operator
(\S\ref{sec:dual}, \S\ref{sec:defect}), and real weight updates can exceed
it in expressivity (\S\ref{sec:ceiling}).
\end{remark}

\subsection{LLMs as approximate posterior simulators}\label{sec:defect}

\begin{fact}[ICL as implicit Bayesian inference~\cite{xie}]\label{fact:xie}
For data generated by a mixture of HMMs (latent concept $\vartheta$),
a sequence model fit to the pretraining distribution performs implicit
Bayesian inference at prompt time: as in-context evidence grows, its
predictive converges to the Bayes-optimal predictor
$\int p(y \mid x, \vartheta)\, d\post(\vartheta \mid \mathrm{prompt})$,
even under distribution mismatch between prompts and pretraining data.
\end{fact}

\begin{definition}[Simulation defect]\label{def:defect}
For a concrete LLM $q_{\theta_0}$, define its \emph{simulation defect} on a
task family as
\[
  \varepsilon_{\mathrm{sim}}(D)
  \;=\; \sup_{x}\;
  d_{\mathrm{TV}}\Bigl( q_{\theta_0}(\cdot \mid x, D),\;
  p_{\mathrm{ctx}}(\cdot \mid x) \Bigr),
\]
the distance between what the model does with context $D$ and what ideal
conditioning would do. All equivalence claims in this document degrade by at
most $\varepsilon_{\mathrm{sim}}$; it plays the role the instrument-noise
$\sigma$ plays in the companion document --- \emph{an empirical constant to
be estimated, not assumed zero}. Documented deviations of real ICL from
idealized learning~\cite{shen} are measurements of
$\varepsilon_{\mathrm{sim}} > 0$, not refutations of the frame.
\end{definition}

%======================================================================
\section{Exact equivalence in the dual regime}\label{sec:dual}
%======================================================================

The Bayesian frame says conditioning and training agree at the level of
predictive distributions. Remarkably, for linear attention the agreement
holds at the level of \emph{mechanism}: stored context consulted by
attention and a weight update computed from the same data are the same
algebraic object.

\begin{theorem}[The dual form: a GD-trained linear layer \emph{is} attention
over its training data~\cite{irie}]\label{thm:dualform}
Let a linear layer $y = W x$, $W \in \R^{m \times n}$, be trained by $T$
steps of gradient descent from $W_0$ on examples $x_1, \dots, x_T \in \R^n$,
producing error signals $\sigma_i \in \R^m$ (the learning-rate-scaled output
deltas $\sigma_i = -\eta_i\, \nabla_{y} \ell_i$). Then for every query
$x \in \R^n$:
\[
  W_T\, x
  \;=\; W_0\, x \;+\; \sum_{i=1}^{T} \bigl( x_i^{\top} x \bigr)\, \sigma_i
  \;=\; W_0\, x \;+\;
  \mathrm{LinAttn}\bigl( \text{query } x;\ \text{keys } \{x_i\},\
  \text{values } \{\sigma_i\} \bigr).
\]
That is, the trained network is \emph{exactly} the initial network plus
unnormalized linear attention over the training set.
\end{theorem}

\begin{proof}
Each GD step adds a rank-one update: $W_{i} = W_{i-1} + \sigma_i x_i^{\top}$.
Summing, $W_T = W_0 + \sum_{i=1}^T \sigma_i x_i^{\top}$. Applying to $x$ and
using $(\sigma_i x_i^{\top}) x = (x_i^{\top} x)\, \sigma_i$ gives the
expression; the right-hand sum is by definition linear attention with keys
$x_i$, values $\sigma_i$, and similarity $\langle \cdot, \cdot \rangle$.
\end{proof}

\begin{remark}
Theorem~\ref{thm:dualform} is an identity, not an approximation. It licenses
the central reading of this document: \textbf{an item retrieved into context
and attended over is, in the linear regime, literally the weight update that
training on it would have produced} --- applied transiently, at inference
time, and removed afterwards. ``Adding context'' and ``taking a gradient
step'' differ in \emph{persistence and coordinates}, not in kind.
\end{remark}

\begin{fact}[Transformers implement learning algorithms in context]
\label{fact:icl-gd}
(i) There exist transformer weights under which one linear self-attention
layer applied to a context of $(x_i, y_i)$ pairs computes the prediction of
one gradient-descent step on the in-context least-squares loss; stacking $L$
layers yields $L$ steps, and \emph{trained} linear-attention models converge
to this construction on toy regression~\cite{vonoswald}. (ii) Transformers
trained on in-context regression learn to match gradient descent, ridge
regression, or exact least squares depending on depth~\cite{akyurek}.
(iii) The global minima of the in-context linear-regression training
objective implement \emph{preconditioned} gradient descent~\cite{ahn}.
(iv) Real-scale evidence: a context induces a compact, transplantable
activation-space \emph{task vector} that reproduces much of the context's
effect --- context compiling to a transient low-rank functional
patch~\cite{taskvectors,functionvectors}. (v) Caveat: pretrained LLM ICL
deviates measurably from exact GD equivalence outside the linear
regime~\cite{shen}; in the present framework these deviations are the
simulation defect of Definition~\ref{def:defect}.
\end{fact}

%======================================================================
\section{The system as a semi-parametric learner}\label{sec:semiparam}
%======================================================================

\begin{definition}[Semi-parametric policy class]\label{def:semiparam}
Fix a frozen base model $\theta_0$. The system's policy class is
\[
  \cF \;=\; \Bigl\{\, f_{L, R} \;:\;
  f_{L,R}(x) = q_{\theta_0}\bigl( \cdot \mid x,\; R(x, L) \bigr)
  \;\Bigm|\; L \in \cL,\ R \in \cR \,\Bigr\},
\]
where $L$ is the insight library (the \emph{nonparametric component}: it
grows with experience, has no fixed dimension, and stores raw knowledge
atoms) and $R$ is the retrieval policy (bounded context budget, as in the
companion document). \emph{Training} is movement in $\cL$ via the governed
operators (insert after quarantine, retire, split, merge); $\theta_0$ never
moves. By Theorem~\ref{thm:dualform} and Fact~\ref{fact:icl-gd}, each
retrieved insight acts as a transient, rank-limited functional update
applied conditionally on the input --- i.e., $\cF$ is a frozen backbone
equipped with a learned, input-gated bank of adapters, stored in the dual
(text) coordinates. The lineage of frozen-core-plus-growing-memory models
(kNN-LM and successors~\cite{knnlm}) is the same architecture one level
down.
\end{definition}

\begin{remark}[``This is not RL,'' upgraded]
The design document's slogan (``no gradient, no policy, no reward model'')
is rhetorically true and semantically misleading. The correct statement:
\emph{this is training} --- of the nonparametric component of
Definition~\ref{def:semiparam}, by discrete, tested updates, in coordinates
where each update is legible. The system's reward signal (the grader) is a
loss; its add/retire operators are steps; its capability curve is a training
curve. What it is \emph{not} is parametric gradient training.
\end{remark}

%======================================================================
\section{The translation dictionary}\label{sec:dictionary}
%======================================================================

The Equivalence Principle makes the following more than analogy: each row
pairs the two coordinate representations of one operation on the predictive
distribution.

\begin{center}
\small
\begin{tabular}{@{}p{0.40\textwidth}p{0.52\textwidth}@{}}
\toprule
\textbf{Weight training} & \textbf{Library learning (the dual)} \\
\midrule
gradient step on a batch &
validated insight insertion (reflector counterfactual) \\
gradient oracle (backprop) &
abduction over traces: a \emph{semantic} gradient, computed by a reasoning
engine with priors \\
step accepted unconditionally &
step accepted only past a level-$\alpha$ anytime-valid gate
(\texttt{theory.tex} \S5) \\
learning rate / trust region &
quarantine strictness $\alpha$, batch size cap \\
early stopping on validation loss &
validation-based reversion (promote-or-revert) \\
regularization; pruning &
generalization lint (prior); retirement / dormancy \\
capacity control (parameter count) &
active cap $k$; streaming-submodular tournament \\
conditional computation / MoE gating &
retrieval policy $R$ (input-dependent adapter selection) \\
catastrophic forgetting (destructive interference in weight space) &
\emph{absent by construction} (append-only store); interference reappears as
routing degradation (shadowing) --- measurable and governable \\
per-example credit: influence functions (approximate, expensive, post hoc) &
per-insight credit: event ledger $+$ randomized ablation (exact by design;
\texttt{theory.tex} \S6) \\
curriculum (data ordering) &
episode/target selection; question-budget annealing \\
training curve (loss vs.\ steps) &
capability curve (ability vs.\ episodes; IRT-anchored) \\
overfitting to the training set &
target memorization (overfitting detector; cross-target recurrence as the
generalization test) \\
checkpoint; learned object dies with it &
library: portable across base-model upgrades \\
distillation into a student &
(future) consolidation of the validated library into weights --- the same
information changing coordinates \\
\bottomrule
\end{tabular}
\end{center}

\begin{remark}[Two rows are improvements, not translations]
In the dual coordinates, two operations are \emph{strictly better} than
their weight-space counterparts: (i) per-update credit is exact (a ledger
row and a designed experiment) where influence functions are heuristic; and
(ii) forgetting is structurally impossible, with its failure energy
displaced into routing, where the companion document's telemetry
(entropy, shadowing margins) can see it. These two rows are the system's
genuine comparative advantages and should anchor the paper's claims.
\end{remark}

%======================================================================
\section{The ceiling: what context optimization cannot do}\label{sec:ceiling}
%======================================================================

\begin{fact}[Elicitation, not creation~\cite{petrov}]\label{fact:petrov}
Context/prefix conditioning of a frozen transformer acts as a bias on the
attention mechanism: it cannot alter the relative attention pattern over
content tokens. Consequently prompting can \emph{elicit} skills represented
by the pretrained model but cannot \emph{create} computational patterns
absent from it. (Prompt tuning approaches fine-tuning quality empirically as
scale grows~\cite{lester}, consistent with large models representing more
elicitable skills, not with the ceiling being void.)
\end{fact}

\begin{proposition}[Context optimization is bounded by weight
optimization]\label{prop:ceiling}
Let $J(f)$ be any performance functional on policies, and suppose the weight
class is \emph{emulation-complete}: for every context policy $(L, R)$ there
exist weights $\theta'$ with
$q_{\theta'}(\cdot \mid x) = q_{\theta_0}(\cdot \mid x, R(x, L))$ for all
$x$. Then
\[
  J^{\ast}_{\mathrm{ctx}}
  \;\coloneqq\; \sup_{L, R}\, J\bigl(f_{L,R}\bigr)
  \;\;\le\;\;
  \sup_{\theta'}\, J\bigl(q_{\theta'}\bigr)
  \;\eqqcolon\; J^{\ast}_{\mathrm{wt}} .
\]
By Fact~\ref{fact:petrov} the inclusion is in general strict.
\end{proposition}

\begin{proof}
Under emulation-completeness, $\{f_{L,R}\} \subseteq \{q_{\theta'}\}$ as
sets of policies; a supremum over a subset is no larger.
\end{proof}

\begin{corollary}[Capability decomposition and the saturation--jump law]
\label{cor:saturation}
Write the attainable objective as $J^\ast_{\mathrm{ctx}}(\theta_0)$ ---
a function of the frozen base model. Library learning moves $J$ toward a
\emph{base-model-dependent frontier} and can go no further; upgrading
$\theta_0$ lifts the frontier while the library (text, model-agnostic up to
re-validation) carries over. The capability curve should therefore show
\emph{saturation toward an asymptote within a base model, and a discrete
jump with library carryover across base-model upgrades}. This is the formal
content of the scaffold-vs-model ability decomposition the measurement layer
estimates psychometrically~\cite{agentpsych}.
\end{corollary}

\begin{remark}[Why the concession is affordable here]
The ceiling binds on \emph{capability} (computational patterns); the
system's accumulated knowledge is overwhelmingly \emph{procedural and
factual} (conventions, pitfalls, workflows, elicitation tactics) --- squarely
inside the elicitable class. The design constraints (no weight access,
subscription budget, auditability, surgical rollback) make weight training
unavailable regardless; Proposition~\ref{prop:ceiling} then prices what is
being given up, and Corollary~\ref{cor:saturation} says how to detect when
the price starts to bind: stalling $p_\Delta$ (companion document,
Thm.~9.2) \emph{with} failure clusters concentrating in
capability-shaped failure kinds.
\end{remark}

%======================================================================
\section{The regime argument: why this does not collapse into poor man's
SGD}\label{sec:regime}
%======================================================================

If conditioning and training are the same operation, is a context-trained
system simply a worse-optimized model? The answer is a regime comparison,
not a verdict.

\begin{definition}[Update regimes]
Call a learning process \emph{averaging-controlled} if it takes many small
updates of low individual semantic content, with quality enforced by
aggregation (SGD: noise averages out; CLT-style analyses apply). Call it
\emph{selection-controlled} if it takes few large updates of high individual
semantic content, with quality enforced by per-update acceptance testing
(this system: each update passes a level-$\alpha$ gate; drift/occupation
analyses apply, cf.\ companion Thm.~9.2).
\end{definition}

\begin{remark}[Where each regime dominates]
Per-update testing costs
$\Theta\bigl(\tfrac{\sigma^2}{\Delta^2}\log\tfrac1\alpha\bigr)$ evaluations
(companion Thm.~5.1/Cor.~5.7); amortizing it requires updates that are
\emph{few and individually valuable} --- exactly what an abductive proposal
engine with world knowledge produces (each ``gradient'' is computed by
reasoning over a failure trace, not by differentiating; the observed
order-of-magnitude rollout efficiencies of reflective prompt evolution over
RL-style search are this effect). Conversely, averaging control needs
update volume that no quota-bound, per-update-tested process can match;
at data scale, parametric training wins. The system's data regime ---
a handful of validated insights per episode --- is the selection-controlled
regime. \textbf{The collapse claim fails not because the equivalence fails,
but because ``poor man's'' presumes the averaging regime, and the system
does not operate there.} What survives of the collapse worry is precisely
Proposition~\ref{prop:ceiling} (the expressivity ceiling) and
Definition~\ref{def:defect} (the simulation defect) --- both quantified,
both monitored.
\end{remark}

\begin{remark}[The distillation corollary]
By the Equivalence Principle, a validated library is simultaneously a
curated dataset: every promoted insight is a (context, behavior) training
pair with a measured causal effect. Consolidation into weights --- when data
scale or the ceiling warrants it --- is the same information changing
coordinates, and the gates' statistics transfer as data-quality
certificates. ``Library learning now, weight learning later'' is therefore
not a pivot but a staged use of one object.
\end{remark}

%======================================================================
\section{Falsifiable predictions}\label{sec:predictions}
%======================================================================

The reframe earns its keep by predicting things the system can test.

\begin{prediction}[Saturation--jump]\label{pred:jump}
Within a fixed base model, the IRT ability curve saturates toward an
asymptote as the library grows; upgrading the base model with the library
held fixed produces a discontinuous ability jump exceeding the recent
within-model trend; library effects (per-insight ablation values) carry over
across the upgrade with mostly preserved sign. (From
Corollary~\ref{cor:saturation}.)
\end{prediction}

\begin{prediction}[Elicitable-failure concentration]\label{pred:elicit}
Library gains concentrate on failure kinds that encode missing
\emph{knowledge or procedure} (convention violations, missed elicitation,
known-pitfall bugs) and are weak on failure kinds that encode missing
\emph{computation} (deep multi-step reasoning shortfalls). The typed failure
taxonomy makes this measurable as a per-failure-kind treatment-effect
profile. (From Fact~\ref{fact:petrov}.)
\end{prediction}

\begin{prediction}[Patch behavior]\label{pred:patch}
A high-value insight's effect behaves like a transient functional patch:
its presence-vs-absence effect on the worker's outputs is reproducible,
approximately localized (affects its scope tag's task cluster, little
else), and approximately additive with unrelated insights --- while
near-duplicate insights interact subadditively (shadowing). Measurable from
the randomized-ablation ledger; subadditivity is the empirical signature of
the submodularity assumption in the companion document
(\texttt{theory.tex} Asm.~8.2). (From Theorem~\ref{thm:dualform}'s
rank-limited-update reading and Fact~\ref{fact:icl-gd}(iv).)
\end{prediction}

\begin{prediction}[Defect scaling]\label{pred:defect}
Where the simulation defect $\varepsilon_{\mathrm{sim}}$ is small
(short, well-structured contexts; tasks resembling pretraining), measured
insight effects track the idealized conditioning account (stable across
paraphrase, order-insensitive); where contexts are long and heterogeneous,
order- and interference-effects grow. This yields a practical design rule
the system already follows for other reasons: \emph{small, structurally
templated, deduplicated insights minimize the defect} --- the structural
schema is, in this frame, defect control.
\end{prediction}

%======================================================================
\section{Ideas retained for further development}\label{sec:revisit}
%======================================================================

Nothing in the earlier line of work is invalidated by the reframe; it
becomes the statistical mechanics of the outer loop around the present
semantics. Retained, with pointers:

\begin{enumerate}[leftmargin=2em]
  \item \textbf{The measurement-limited learning rate}
        (\texttt{theory.tex} Thm.~5.1, Cor.~5.2) --- now reads as: the
        \emph{trainer's} step-validation bandwidth is bounded by instrument
        noise; in training language, the gate is the system's
        learning-rate schedule.
  \item \textbf{Anytime-valid gates, e-value composition, e-BH}
        (\texttt{theory.tex} \S5) --- per-step validation tests; in this
        frame, a provably safe optimizer in the selection-controlled
        regime.
  \item \textbf{Causal credit: randomized masking, Shapley}
        (\texttt{theory.tex} \S6) --- the dual's exact replacement for
        influence functions.
  \item \textbf{Submodular governance of the cap}
        (\texttt{theory.tex} \S7) --- capacity control of the nonparametric
        component; Prediction~\ref{pred:patch} ties its key assumption to
        the patch picture.
  \item \textbf{Drift/occupation convergence} (\texttt{theory.tex}
        Thm.~9.2) --- convergence of selection-controlled training to
        statistical stationarity; $p_\Delta$ decay is the training-loss
        plateau, and \S\ref{sec:ceiling} now distinguishes \emph{two}
        plateau causes (proposal exhaustion vs.\ elicitation ceiling) with
        Prediction~\ref{pred:elicit} as the discriminator.
  \item \textbf{The adversarial instrument and exploitability bound}
        (\texttt{theory.tex} \S10) --- in training language: a
        non-stationary loss that provably resists reward hacking at the
        falsifier's regret rate.
  \item \textbf{The grader-as-loss program} (ideation record
        2026-06-10): STL robustness margins (differentiable --- and now
        doubly motivated: a graded loss is what a future weight-coordinate
        consolidation would train against), bisimulation distance,
        constrained (CMDP) composition, pairwise judging, conformal triage.
  \item \textbf{Impossibility boundary and Assumption Ledger}
        (\texttt{theory.tex} \S\S10--11) --- extended by one row:
        $\varepsilon_{\mathrm{sim}}$ (Definition~\ref{def:defect}),
        estimable via paraphrase/order-invariance probes on frozen replay
        scenarios (Prediction~\ref{pred:defect}).
\end{enumerate}

%======================================================================
\section{Position for the paper}
%======================================================================

The claimable contribution, in one paragraph: \emph{skill-library learning
systems (Voyager, AWM, ACE, GEPA descendants) have been described as prompt
engineering with memory; we give them training-theoretic semantics. Via the
de Finetti equivalence (conditioning $=$ parameter update on predictives)
and the dual-form identity (attention over stored data $=$ a GD-trained
layer), a governed text library is the nonparametric component of a
semi-parametric model, and library curation is training in the dual
coordinates --- in a selection-controlled regime with exact per-update
credit, no destructive forgetting, statistical step-validation, and an
elicitation ceiling that the system's own psychometrics can detect. The
framework yields falsifiable predictions (saturation--jump, elicitable
failure concentration, patch additivity) testable on the system's own
telemetry.} To our knowledge no published work connects this citation chain
to skill libraries with a formal statement; the nearest neighbors treat ICL
theory and memory systems separately.

%======================================================================
\begin{thebibliography}{99}\small

\bibitem{definetti} B.~de Finetti.
La pr\'evision: ses lois logiques, ses sources subjectives.
\emph{Annales de l'Institut Henri Poincar\'e}, 7:1--68, 1937.

\bibitem{bernardosmith} J.~Bernardo and A.~Smith.
\emph{Bayesian Theory}. Wiley, 1994. (Representation theorems, Ch.~4.)

\bibitem{xie} S.~M. Xie, A.~Raghunathan, P.~Liang, T.~Ma.
An explanation of in-context learning as implicit Bayesian inference.
\emph{ICLR}, 2022. arXiv:2111.02080.

\bibitem{irie} K.~Irie, I.~Schlag, R.~Csord\'as, J.~Schmidhuber.
The dual form of neural networks revisited: connecting test-time predictions
to training patterns via spotlights of attention.
\emph{ICML}, 2022. arXiv:2202.05798.

\bibitem{vonoswald} J.~von Oswald, E.~Niklasson, E.~Randazzo,
J.~Sacramento, A.~Mordvintsev, A.~Zhmoginov, M.~Vladymyrov.
Transformers learn in-context by gradient descent.
\emph{ICML}, 2023. arXiv:2212.07677.

\bibitem{akyurek} E.~Aky\"urek, D.~Schuurmans, J.~Andreas, T.~Ma, D.~Zhou.
What learning algorithm is in-context learning? Investigations with linear
models. \emph{ICLR}, 2023. arXiv:2211.15661.

\bibitem{ahn} K.~Ahn, X.~Cheng, H.~Daneshmand, S.~Sra.
Transformers learn to implement preconditioned gradient descent for
in-context learning. \emph{NeurIPS}, 2023. arXiv:2306.00297.

\bibitem{dai} D.~Dai, Y.~Sun, L.~Dong, Y.~Hao, S.~Ma, Z.~Sui, F.~Wei.
Why can GPT learn in-context? Language models implicitly perform gradient
descent as meta-optimizers. \emph{Findings of ACL}, 2023. arXiv:2212.10559.

\bibitem{shen} L.~Shen, A.~Mishra, D.~Khashabi.
Do pretrained transformers really learn in-context by gradient descent?
arXiv:2310.08540, 2023.

\bibitem{taskvectors} R.~Hendel, M.~Geva, A.~Globerson.
In-context learning creates task vectors. \emph{Findings of EMNLP}, 2023.
arXiv:2310.15916.

\bibitem{functionvectors} E.~Todd, M.~Li, A.~Sharma, A.~Mueller,
B.~Wallace, D.~Bau. Function vectors in large language models.
\emph{ICLR}, 2024. arXiv:2310.15213.

\bibitem{petrov} A.~Petrov, P.~Torr, A.~Bibi.
When do prompting and prefix-tuning work? A theory of capabilities and
limitations. \emph{ICLR}, 2024. arXiv:2310.19698.

\bibitem{lester} B.~Lester, R.~Al-Rfou, N.~Constant.
The power of scale for parameter-efficient prompt tuning.
\emph{EMNLP}, 2021. arXiv:2104.08691.

\bibitem{knnlm} U.~Khandelwal, O.~Levy, D.~Jurafsky, L.~Zettlemoyer,
M.~Lewis. Generalization through memorization: nearest neighbor language
models. \emph{ICLR}, 2020. arXiv:1911.00172.

\bibitem{agentpsych} (Agent psychometrics.)
Task-level performance prediction in agentic coding benchmarks:
decomposing model and scaffold ability. arXiv:2604.00594, 2026.

\end{thebibliography}

\end{document}
